\documentclass[10pt,aspectratio=169]{beamer}
% Preámbulo modular para presentaciones Beamer (Modelación Estadística)
% Diseño y paleta institucional del Tecnológico de Monterrey

\documentclass[aspectratio=169,xcolor={dvipsnames,table}]{beamer}

\usepackage[utf8]{inputenc}
\usepackage[spanish,mexico]{babel}
\usepackage{amsmath,amssymb,amsfonts,mathtools}
\usepackage{graphicx}
\usepackage{listings}
\usepackage{booktabs}
\usepackage{tikz}
\usetikzlibrary{matrix,arrows,decorations.pathmorphing,shapes.geometric,positioning}

% Paleta Institucional Tecnológico de Monterrey
\definecolor{TecAzul}{HTML}{0039A6}       % Azul TEC: color primario institucional
\definecolor{TecAzulOscuro}{HTML}{1A2E51} % Azul marino oscuro
\definecolor{TecRojo}{HTML}{EC2661}       % Rojo de acento (paleta miTec)
\definecolor{TecGris}{HTML}{646464}       % Gris neutro para comentarios y subtítulos
\definecolor{TecFondoBloque}{HTML}{F4F6F9}% Fondo suave para bloques

% Configuración de Tema Beamer
\usetheme{Boadilla}
\usecolortheme[named=TecAzulOscuro]{structure}
\setbeamercolor{alerted text}{fg=TecRojo}
\setbeamercolor{palette primary}{bg=TecAzulOscuro,fg=white}
\setbeamercolor{palette secondary}{bg=TecAzul,fg=white}
\setbeamercolor{palette tertiary}{bg=TecAzulOscuro!80!black,fg=white}
\setbeamercolor{title}{fg=TecAzulOscuro,bg=TecFondoBloque}
\setbeamercolor{frametitle}{fg=TecAzulOscuro,bg=TecFondoBloque}

% Bloques personalizados elegantes
\setbeamercolor{block title}{bg=TecAzulOscuro,fg=white}
\setbeamercolor{block body}{bg=TecFondoBloque,fg=black}
\setbeamercolor{block title alerted}{bg=TecRojo,fg=white}
\setbeamercolor{block body alerted}{bg=TecRojo!8,fg=black}
\setbeamercolor{block title example}{bg=TecAzul,fg=white}
\setbeamercolor{block body example}{bg=TecAzul!8,fg=black}

% Redondear bordes y añadir sombra sutil
\setbeamertemplate{blocks}[rounded][shadow=true]
\setbeamertemplate{navigation symbols}{}

% Pie de página limpio con número de diapositiva
\setbeamertemplate{footline}{
  \leavevmode%
  \hbox{%
  \begin{beamercolorbox}[wd=.7\paperwidth,ht=2.25ex,dp=1ex,left]{author in head/foot}%
    \hspace*{2ex}\insertshortauthor~~(\insertshortinstitute) \hfill \insertshorttitle
  \end{beamercolorbox}%
  \begin{beamercolorbox}[wd=.3\paperwidth,ht=2.25ex,dp=1ex,right]{date in head/foot}%
    \insertshortdate{}\hspace*{2em}
    \insertframenumber{} / \inserttotalframenumber\hspace*{2ex}
  \end{beamercolorbox}}%
  \vskip0pt%
}

% Configuración de Listings de Python para visualización pedagógica de código
\lstset{
    language=Python,
    basicstyle=\ttfamily\footnotesize,
    keywordstyle=\color{TecAzulOscuro}\bfseries,
    stringstyle=\color{TecRojo},
    commentstyle=\color{TecGris}\itshape,
    showstringspaces=false,
    breaklines=true,
    frame=lines,
    rulecolor=\color{TecAzulOscuro!30},
    backgroundcolor=\color{TecFondoBloque},
    aboveskip=0.5em,
    belowskip=0.5em,
    literate={á}{{\'a}}1 {é}{{\'e}}1 {í}{{\'i}}1 {ó}{{\'o}}1 {ú}{{\'u}}1
             {Á}{{\'A}}1 {É}{{\'E}}1 {Í}{{\'I}}1 {Ó}{{\'O}}1 {Ú}{{\'U}}1
             {ñ}{{\~n}}1 {Ñ}{{\~N}}1 {¿}{{?`}}1 {¡}{{!`}}1
}

% Metadatos institucionales por defecto
\title[Modelación Estadística]{Modelación Estadística para Ciencia de Datos e Ingeniería}
\author[J. Castillo Colmenares]{Juliho David Castillo Colmenares}
\institute[Tec de Monterrey]{Tecnológico de Monterrey \\ Escuela de Ingeniería y Ciencias}
\date{\today}
% Comandos y macros estadísticas compartidas para las presentaciones Beamer (Metropolis)

% Conjuntos numéricos básicos
\newcommand{\R}{\mathbb{R}}
\newcommand{\N}{\mathbb{N}}
\newcommand{\Q}{\mathbb{Q}}
\newcommand{\Z}{\mathbb{Z}}
\newcommand{\set}[1]{\left\{ #1 \right\}}
\newcommand{\abs}[1]{\left|#1\right|}
\newcommand{\norm}[1]{\left\| #1 \right\|}

% Comandos estadísticos y probabilísticos del curso
\newcommand{\Var}{\operatorname{Var}}
\newcommand{\cov}{\operatorname{Cov}}
\newcommand{\corr}{\rho}
\newcommand{\comb}[2]{\begin{pmatrix} #1 \\ #2 \end{pmatrix}}
\newcommand{\s}{\sigma}
\newcommand{\particion}{\mathfrak{P}}
\newcommand{\card}[1]{\#\left( #1 \right)}
\newcommand{\seq}[3]{\set{#1}_{#2}^{#3}}
\newcommand{\E}{\mathbb{E}}
\newcommand{\Prob}{\mathbb{P}}

% Entornos pedagógicos y cajas en español académico natural
\newenvironment{discusion}{%
  \begin{alertblock}{Pregunta de Discusión}%
}{%
  \end{alertblock}%
}

\newenvironment{reflexion}{%
  \begin{alertblock}{Reflexión para el Aula}%
}{%
  \end{alertblock}%
}

\newenvironment{paradoja}{%
  \begin{alertblock}{Paradoja de Exploración}%
}{%
  \end{alertblock}%
}

\newenvironment{motivacion}{%
  \begin{exampleblock}{Motivación: Ciencia de Datos en la Práctica}%
}{%
  \end{exampleblock}%
}

\newenvironment{retoclase}{%
  \begin{block}{Reto Rápido en Equipo}%
}{%
  \end{block}%
}

\title[Resumen de modelos]{Sección 09.11: Resumen de modelos}\subtitle{Criterios para seleccionar predictores}\author[J. Castillo Colmenares]{Juliho Castillo Colmenares}\institute{Tecnológico de Monterrey}\date{}
\begin{document}
\begin{frame}[plain]\titlepage\end{frame}
\begin{frame}{Hoja de ruta}\small\begin{enumerate}\item Qué significa un buen modelo.\item Métricas de ajuste y complejidad.\item Diagnóstico y selección responsable.\end{enumerate}\end{frame}
\begin{frame}{Fuente y alcance}\small La sección resume criterios para construir un modelo lineal confiable y orienta la selección de variables predictoras.\begin{block}{Pregunta guía}¿Cómo comparar especificaciones sin premiar complejidad innecesaria?\end{block}\end{frame}
\begin{frame}{Criterio general}\small Un buen modelo combina ajuste, parsimonia, estabilidad, interpretabilidad y desempeño fuera de muestra. Ninguna métrica aislada decide por sí sola.\end{frame}
\begin{frame}{$R^2$}\small $R^2$ mide la fracción de variabilidad explicada en la muestra. Nunca disminuye al añadir predictores, incluso si la nueva variable aporta poco.\end{frame}
\begin{frame}{$R^2$ ajustado}\small El $R^2$ ajustado penaliza el número de predictores. Puede disminuir cuando una variable nueva no mejora suficientemente el ajuste relativo a su costo.\end{frame}
\begin{frame}{Valor $p$ de un predictor}\small Un valor $p$ pequeño para el coeficiente aporta evidencia contra una contribución nula bajo el modelo, pero depende de supuestos y de la escala de análisis.\end{frame}
\begin{frame}{Estadístico $F$ global}\small El estadístico $F$ contrasta conjuntamente si el modelo explica variación más allá del intercepto. Su valor $p$ evalúa la especificación global.\end{frame}
\begin{frame}{Error estándar residual}\small El RSE resume la dispersión típica de los residuos en unidades de la respuesta. Menor RSE indica menor error dentro de la muestra, no necesariamente mejor generalización.\end{frame}
\begin{frame}{Factor de inflación de varianza}\small El VIF detecta cuánto se infla la varianza de un coeficiente por colinealidad. Valores altos advierten que los efectos parciales son inestables.\end{frame}
\begin{frame}{Complejidad}\small Añadir predictores siempre puede mejorar $R^2$, pero aumenta grados de libertad consumidos, dificultad de interpretación y riesgo de sobreajuste.\end{frame}
\begin{frame}{Selección de variables}\small La selección debe comenzar con la pregunta científica, disponibilidad de datos y plausibilidad de los predictores; las métricas ayudan a comparar, no sustituyen el criterio.\end{frame}
\begin{frame}{Procedimiento I}\small\begin{enumerate}\item Definir respuesta y candidatos.\item Ajustar un modelo base.\item Revisar $R^2$, RSE y $F$.\end{enumerate}\end{frame}
\begin{frame}{Procedimiento II}\small\begin{enumerate}\setcounter{enumi}{3}\item Comparar $R^2$ ajustado.\item Revisar valores $p$ con cautela.\item Diagnosticar VIF y residuos.\end{enumerate}\end{frame}
\begin{frame}{Aplicación: añadir predictor}\small Al añadir una variable, $R^2$ sube automáticamente. La decisión exige comprobar si también mejora $R^2$ ajustado, reduce RSE y mantiene sentido sustantivo.\end{frame}
\begin{frame}{Aplicación: $F$ global}\small Un aumento del estadístico $F$ y una disminución de su valor $p$ pueden indicar mejora global, pero deben interpretarse con los grados de libertad y la validación.\end{frame}
\begin{frame}{Aplicación: colinealidad}\small Si dos predictores contienen información casi redundante, el VIF puede ser alto aunque el ajuste global sea excelente. Considerar eliminar, combinar o regularizar variables.\end{frame}
\begin{frame}{Aplicación: validación}\small Comparar error de entrenamiento y de validación. Un modelo ligeramente peor en muestra puede predecir mejor fuera de muestra si es más estable.\end{frame}
\begin{frame}{Interpretación}\small Los criterios son complementarios: ajuste ($R^2$), penalización ($R^2$ ajustado), evidencia ($p$, $F$), escala de error (RSE) y estabilidad (VIF).\end{frame}
\begin{frame}{Reporte}\small Documente variables candidatas, especificaciones comparadas, métricas, diagnósticos, criterio de selección y limitaciones.\end{frame}
\begin{frame}{Síntesis}\small\begin{itemize}\item No optimizar una sola cifra.\item Penalizar complejidad.\item Revisar colinealidad y residuos.\item Validar fuera de muestra.\end{itemize}\end{frame}
\begin{frame}{Conexión con el capítulo}\small Este resumen integra regresión, validación, variables categóricas y transformaciones en una decisión reproducible de modelado.\end{frame}
\end{document}
